% page 633 of the facsimile (image p-632 of the scan)
% block 170.2 x 242.0 mm ; left margin 31.8 mm
\pgc{10.160mm}{0.000mm}{
\l{}
\l{\cel{56}625}
\l{}
\l{\sou{yako}. (\sou{zool.}) Bufalo kun kaval-kaudo, devenanta de Tibet. -}
\l{\cel{4}L.\cel{1}\sou{poephagus}. - F.}
\l{}
\l{\sou{yakto}.\cel{2}Exkurso-navo. - DEFRS.}
\l{}
\l{\sou{yaro}. I. Duro di un jiro da la Terglobo cirkum Suno (365 dii,}
\l{\cel{4}5 hori, 48 minuti, e 50 sekundi), egardata kom unajo di la}
\l{\cel{4}mezuro di tempo. - II. Ta periodo determinata per ula quanto}
\l{\cel{4}de fakti qui eventis dum olca. - III. Periodo de dek-e-du}
\l{\cel{4}monati, egardata relate to quo eventas dum li, sen egardo}
\l{\cel{4}pri la epoko kande li komencas. -IV. Singla periodo de dek-}
\l{\cel{4}e-du monati, kontita de la nasko-dio di persono. - DE.}
\l{}
\l{\sou{yardo}. (\sou{navig.})\cel{2}Ligno-peco, dinigita ye lua extremaji, lokizita}
\l{\cel{4}horizontale ye altesi diversa, sur la masti di navo, por}
\l{\cel{4}subtenar la veli desvolvita o volvita. - E.}
\l{}
\l{\sou{yatagano}. Sabro poniardo-forma, kun lamo kelke kurva proxim la}
\l{\cel{4}pinto - armo Araba e Turka. - DEFIS.}
\l{}
\l{ye. Prepoziciono di qua la senco ne esas determinita, quan onu}
\l{\cel{4}uzas nur en la kazi kande nula propoziciono altra postul-}
\l{\cel{4}\sou{esas da la senco}. (Lu indikas nome la loko o la dato exakta}
\l{\cel{4}di evento, di fakto, la parto koncernata di la korpo: \textquotedbl{}ye}
\l{\cel{4}la angulo di la strado\textquotedbl{}, \textquotedbl{}ye la dekesma kilometro\textquotedbl{}, \textquotedbl{}ye}
\l{\cel{4}dimezo\textquotedbl{}, \textquotedbl{}ye la lasta foyo\textquotedbl{}, \textquotedbl{}me doloras ye la kapo\textquotedbl{}, \textquotedbl{}ilu}
\l{\cel{4}prenis elu ye la tayo\textquotedbl{}, \textquotedbl{}il kaptis la kavalo per lazo ye}
\l{\cel{4}la kolo\textquotedbl{}, \textquotedbl{}plenigar ulo ye ... per kuliero, spado, e c.\textquotedbl{})}
\l{}
\l{\sou{yen}. Prepoziciono qua igas atencar la persono o la kozo pri qua}
\l{\cel{4}onu quik parolos, o to quo quik eventos. (\sou{yen}, avan nomo ;}
\l{\cel{4}\sou{yen}\cel{1}ke, avan propoziciono).}
\l{}
\l{\sou{yes}.\cel{2}Adverbo qua equivalas propoziciono per qua onu respondas}
\l{\cel{4}afirme a questiono (expresita o tacita). (\sou{yes}\cel{1}ya = +si).- E.}
\l{}
\l{\sou{yiterbo.}\cel{2}(\sou{kemio)}\cel{2}Substanco ne-izolita, di qua onu konocas nur}
\l{\cel{4}la oxido, e qua ne esus altro kam mixuro de lutecio e di neo-}
\l{\cel{4}yiterbio. -\cel{1}\sou{Simbolo kemiala}\cel{1}: Yb.}
\l{}
\l{\sou{yitro}. (\sou{kemio)}.\cel{2}Metalo skarsa, qua akompanas cerio en la maxim}
\l{\cel{4}multa de lua erci. Atom-pezo : 88,92. -\cel{1}\sou{Simb. kemiala}: Y.}
\l{}
\l{\sou{yodlar}. (\sou{netrans.})\cel{2}Kantar ta kansono montanala en qua onu gutur-}
\l{\cel{4}agas stroke por pasar de la voco \textquotedbl{}pektorala\textquotedbl{} a la voco \textquotedbl{}kapala\textquotedbl{}}
\l{\cel{4}(voco di qua la rezono eventas en la kavaji supra od infra,}
\l{\cel{4}rispektive, di la organaro vocala), di qua la rezultajo esas}
\l{\cel{4}quaza rukulo tre stranja. - DEF.}
\l{}
\l{\sou{yolo}. (\sou{navig.})\cel{2}Mikra kanoto, streta e lejera, qua avancas rapide.}
\l{\cel{4}- DEFIR.}
\l{}
\l{\sou{yufto}. Ledro traktita per birko-oleo. - DR.}
}
